% Module 4 section file. Included by ../module4_alfven_continuum_eigenmodes.tex.
\section{Construction of the reduced coupled-harmonic benchmark}
\label{sec:benchmark_construction}

\subsection{Design philosophy}
A verification benchmark should be complicated enough to exercise the intended algorithms but simple enough that its scaling, symmetry, limits, and expected behavior can be derived independently. The Module 4 benchmark therefore does not import a numerical equilibrium. It prescribes smooth analytic profiles and constructs a real two-harmonic operator with controlled coupling.

The design requirements are:
\begin{enumerate}[leftmargin=1.6em]
\item all coefficients have explicit SI dimensions or are explicitly dimensionless;
\item the uncoupled limit is analytically known;
\item the local coupling opens a finite avoided crossing;
\item the global operator is self-adjoint and positive for the baseline case;
\item the Alfv\'en scaling with $B_0$ and density is exact;
\item the finite-volume matrix is real and symmetric;
\item Python and MATLAB can implement the same mathematical contract independently.
\end{enumerate}

\subsection{Normalized radial coordinate}
The coordinate $s\in[0,1]$ is described in the case metadata as a normalized toroidal-flux radius. In the benchmark equations it is a dimensionless monotonic surface label. It should not be assumed identical to geometric minor radius $r/a$ or to the square root of normalized toroidal flux unless a future imported equilibrium defines that mapping explicitly.

The practical interpretation is:
\begin{equation}
s=0\quad\text{axis surrogate},
\qquad
s=1\quad\text{edge surrogate}.
\end{equation}
All radial derivatives in the implemented model are derivatives with respect to this dimensionless coordinate.

\subsection{Mass-density profile}
The prescribed mass density is
\begin{equation}
\boxed{
\rho(s)=\rho_0\left[1-(1-f_\rho)s^{a_\rho}\right].}
\label{eq:benchmark_density}
\end{equation}
Here:
\begin{itemize}[leftmargin=1.6em]
\item $\rho_0>0$ is the axis density;
\item $f_\rho=\rho(1)/\rho_0$ is the edge-to-axis density fraction;
\item $a_\rho>0$ controls profile shape.
\end{itemize}
At $s=0$, $\rho=\rho_0$. At $s=1$, $\rho=f_\rho\rho_0$. If $0<f_\rho\le1$, density remains positive. Its derivative is
\begin{equation}
\rho'(s)=-\rho_0(1-f_\rho)a_\rho s^{a_\rho-1}.
\end{equation}
For the baseline $a_\rho=2$ and $f_\rho=0.2$, density decreases smoothly toward the edge, so $v_A$ increases outward.

\subsection{Rotational-transform profile}
The prescribed transform is
\begin{equation}
\boxed{
\iota(s)=\iota_0+(\iota_1-\iota_0)s^{a_\iota}.}
\label{eq:benchmark_iota}
\end{equation}
The parameters $\iota_0$ and $\iota_1$ are the axis and edge values, and $a_\iota>0$ controls radial shape. Its derivative is
\begin{equation}
\iota'(s)=(\iota_1-\iota_0)a_\iota s^{a_\iota-1}.
\end{equation}
For the baseline $a_\iota=1$, the transform varies linearly.

If a target transform $\iota_*$ lies between $\iota_0$ and $\iota_1$, the inverse profile is
\begin{equation}
\boxed{
s(\iota_*)=
\left(\frac{\iota_*-\iota_0}{\iota_1-\iota_0}\right)^{1/a_\iota}.}
\label{eq:iota_inverse}
\end{equation}
The code clips the fraction to $[0,1]$ as a defensive operation, but a physically meaningful intended crossing should lie within the prescribed transform range.

\subsection{Alfv\'en speed and frequency scale}
Using a constant reference magnetic-field magnitude $B_0$,
\begin{equation}
\boxed{v_A(s)=\frac{B_0}{\sqrt{\mu_0\rho(s)}}.}
\label{eq:benchmark_va}
\end{equation}
Define
\begin{equation}
\omega_0(s)=\frac{v_A(s)}{R_0}.
\label{eq:benchmark_omega_scale}
\end{equation}
The quantity $\omega_0$ has units $\mathrm{s^{-1}}$. Every term in the benchmark operator is proportional to $\omega_0^2$, ensuring exact frequency scaling with $B_0$ and $\rho^{-1/2}$.

\subsection{Uncoupled branches}
For fixed $n$ and neighboring poloidal harmonics $m$ and $m+1$,
\begin{align}
\omega_m(s)&=|n-m\iota(s)|\omega_0(s),
\label{eq:benchmark_branch_m}\\
\omega_{m+1}(s)&=|n-(m+1)\iota(s)|\omega_0(s).
\label{eq:benchmark_branch_mp1}
\end{align}
These are local positive angular frequencies. The corresponding signed parallel wave numbers retain propagation direction, but only their squares enter the real symmetric benchmark.

\subsection{Analytic crossing location}
The TAE-like equal-magnitude opposite-sign condition gives Eq.~\eqref{eq:tae_crossing_iota}:
\begin{equation}
\iota_c=\frac{2n}{2m+1}.
\end{equation}
Substitution into Eq.~\eqref{eq:iota_inverse} gives
\begin{equation}
\boxed{
s_c=
\left[
\frac{2n/(2m+1)-\iota_0}{\iota_1-\iota_0}
\right]^{1/a_\iota}.}
\label{eq:benchmark_crossing_s}
\end{equation}
This formula is valid when the bracketed fraction lies in $[0,1]$.

\subsection{Localized coupling}
The benchmark coupling is
\begin{equation}
\boxed{
C(s)=\epsilon_c\omega_m(s)\omega_{m+1}(s)
\exp\left[-\left(\frac{s-s_c}{w_c}\right)^2\right].}
\label{eq:benchmark_coupling}
\end{equation}
The dimensionless parameter $\epsilon_c\in[0,1)$ controls coupling strength, and $w_c>0$ controls radial width. The Gaussian is centered at the analytic crossing because that is where neighboring harmonics are expected to mix most strongly in the illustrative model.

The coupling is a surrogate. It is not derived from an equilibrium Fourier coefficient. Its value is positive in the baseline. A negative coupling would change the relative sign of mixed eigenvectors but not the two local eigenvalues because Eq.~\eqref{eq:two_mode_eigenvalues} depends on $C^2$.

\subsection{Local coupled branches}
At each grid coordinate, define
\begin{equation}
\mathbf H_{\mathrm{loc}}(s)=
\begin{bmatrix}
\omega_m^2(s)&C(s)\\
C(s)&\omega_{m+1}^2(s)
\end{bmatrix}.
\end{equation}
The coupled local frequencies are
\begin{equation}
\boxed{
\Omega_\pm^2(s)=
\frac{\omega_m^2+\omega_{m+1}^2}{2}
\pm
\sqrt{
\left(\frac{\omega_m^2-\omega_{m+1}^2}{2}\right)^2+C^2
}.}
\label{eq:benchmark_local_branches}
\end{equation}
The code uses $\sqrt{\max(\Omega_\pm^2,0)}$ to protect against tiny negative roundoff before taking a square root. For valid baseline parameters, the analytic positivity condition should already keep $\Omega_\pm^2>0$.

\subsection{Radial stiffness}
The radial coefficient is
\begin{equation}
\boxed{
D(s)=\epsilon_r\left(\frac{v_A(s)}{R_0}\right)^2,}
\label{eq:benchmark_D}
\end{equation}
where $\epsilon_r>0$ is dimensionless. Since $s$ is dimensionless, $d/ds$ carries no physical length unit, so $D$ must have units of squared angular frequency. This is why a factor $a^{-2}$ does not appear in the implemented equation. A model using geometric radius $r$ instead would require a coefficient with units of speed squared.

\subsection{Dimensionless form}
Choose an axis frequency scale
\begin{equation}
\omega_{00}=\frac{B_0}{R_0\sqrt{\mu_0\rho_0}}.
\end{equation}
Define
\begin{equation}
\widehat\omega=\frac{\omega}{\omega_{00}},
\qquad
\widehat\rho=\frac{\rho}{\rho_0},
\qquad
\widehat D=\frac{D}{\omega_{00}^2}.
\end{equation}
Then
\begin{align}
\frac{v_A/R_0}{\omega_{00}}&=\widehat\rho^{-1/2},\\
\widehat\omega_m&=|n-m\iota|\widehat\rho^{-1/2},\\
\widehat D&=\epsilon_r\widehat\rho^{-1}.
\end{align}
The entire dimensionless spectrum depends only on profile shapes, mode numbers, $\epsilon_c$, $w_c$, $\epsilon_r$, and boundary conditions. Multiplying $B_0$ or $R_0^{-1}$ rescales all dimensional frequencies without changing dimensionless eigenvectors.

\subsection{Exact scaling relations}
If $B_0\rightarrow\alpha B_0$ with density unchanged, then
\begin{equation}
v_A\rightarrow\alpha v_A,
\quad
\omega_m^2,C,D\rightarrow\alpha^2(\omega_m^2,C,D),
\quad
\omega_j\rightarrow\alpha\omega_j.
\label{eq:B_scaling}
\end{equation}
If $\rho(s)\rightarrow\beta\rho(s)$ uniformly, then
\begin{equation}
v_A\rightarrow\beta^{-1/2}v_A,
\quad
\omega_j\rightarrow\beta^{-1/2}\omega_j.
\label{eq:rho_scaling}
\end{equation}
If $R_0\rightarrow\chi R_0$ with all dimensionless profiles unchanged, then
\begin{equation}
\omega_j\rightarrow\chi^{-1}\omega_j.
\label{eq:R_scaling}
\end{equation}
These are stronger tests than merely reproducing one numerical table because they interrogate every matrix block simultaneously.
